\documentclass[11pt,a4paper]{article}

\usepackage[margin=1in]{geometry}
\usepackage{amsmath, amssymb}
\usepackage{booktabs}
\usepackage{tabularx}
\usepackage{array}
\usepackage{hyperref}
\usepackage{enumitem}
\usepackage{xcolor}
\usepackage{listings}
\usepackage{graphicx}
\usepackage{tcolorbox}

% Code listings
\definecolor{codebg}{rgb}{0.96,0.96,0.96}
\definecolor{codecomment}{rgb}{0.25,0.5,0.35}
\definecolor{codekw}{rgb}{0.6,0.0,0.5}
\lstset{
  basicstyle=\ttfamily\footnotesize,
  backgroundcolor=\color{codebg},
  commentstyle=\color{codecomment}\itshape,
  keywordstyle=\color{codekw}\bfseries,
  breaklines=true,
  numbers=none,
  frame=single,
  framerule=0pt,
  framesep=4pt,
  showstringspaces=false,
  xleftmargin=0pt,
  xrightmargin=0pt,
}

\hypersetup{
  colorlinks=true,
  linkcolor=blue!50!black,
  urlcolor=blue!50!black,
}

\title{
  \vspace{-1cm}
  \textbf{Multi-Agent Loan Default Prediction System}\\
  \large System Design Document
}
\author{Architecture overview, agent responsibilities, and design rationale}
\date{}

\begin{document}
\maketitle
\vspace{-1cm}

\tableofcontents
\vspace{0.5cm}

%==============================================================
\section{Overview}
\label{sec:overview}

The system predicts loan default for Grade C--G LendingClub-style loans by
fusing a numeric XGBoost baseline with text-derived risk signals extracted
from the borrower's free-text description. Five LLM-based agents and one
learned per-row gating function collaborate to decide \emph{when} and
\emph{how strongly} to trust the text.

\subsection{Pipeline at a glance}

\begin{center}
\begin{tabular}{rl}
\textbf{1.} & \textbf{Green --- XGBoost Baseline} (numeric model, always fires) \\
\textbf{2.} & \textbf{Text Analyst} (LLM, always fires if description present) \\
\textbf{3.} & \textbf{Tier-A learned gate} (per-row text-weight, always fires) \\
\textbf{4.} & \textbf{Feature Strategist} (LLM, conditionally fires) \\
\textbf{5.} & \textbf{Subgroup Advocate} (LLM, fires only on Strategist deviation) \\
\textbf{6.} & \textbf{Green --- Arbitration} (LLM, fires only on Advocate veto) \\
\textbf{7.} & Fusion (logit-space math) \\
\textbf{8.} & \textbf{Green --- Final Authority} (hard-rule flip gate) \\
\textbf{9.} & \textbf{Green --- LLM Review} (LLM, fires only on flip block) \\
\textbf{10.} & \textbf{Reporter} (LLM, generates explanation)
\end{tabular}
\end{center}

\noindent
The Green Agent owns multiple roles: baseline scoring, fusion validation,
arbitration on veto, the Tier-A gate training, the final-authority flip gate,
and the terminal LLM review. The previously-separate Arbitrator agent has been
merged into Green (see \S\ref{sec:design-arbitrator}).

\subsection{Fusion formula}

All probabilities are fused in logit space:
\begin{align*}
\mathrm{fused\_logit} &= \mathrm{logit}(p_{\mathrm{base}})
                       + w \cdot c \cdot \bigl(\mathrm{logit}(p_{\mathrm{text}}) - \mathrm{logit}(0.5)\bigr) \\
\mathrm{fused\_pred}  &= \sigma(\mathrm{fused\_logit})
\end{align*}
where $p_{\mathrm{base}}$ is the XGBoost baseline probability, $p_{\mathrm{text}}$
is the Text Analyst's risk score, $c$ is its confidence, and $w$ is the text weight
produced by the Tier-A gate (or by the static 5-scalar policy if no gate is loaded).

%==============================================================
\section{Agent responsibilities}
\label{sec:agents}

\subsection{Text Analyst (LLM, White Agent 1)}

\textbf{Input:} borrower description text (truncated to 500 chars).

\noindent
\textbf{Output:} \texttt{text\_risk\_score} $\in[0.05, 0.95]$ and
\texttt{confidence} $\in[0.4, 0.95]$, plus a list of identified
\texttt{risk\_signals} and \texttt{protective\_signals}, and a short
\texttt{reasoning} sentence.

\noindent
\textbf{Scoring rubric (in the LLM prompt):} start at 0.45, then add or
subtract:
\begin{itemize}[itemsep=0pt]
    \item $-0.15$ stable employment or specific job title
    \item $-0.10$ concrete repayment plan
    \item $-0.10$ ``no missed payments''
    \item $-0.08$ specific income amount
    \item $-0.05$ single clear purpose
    \item $+0.15$ urgent or desperate language
    \item $+0.12$ behind on payments
    \item $+0.10$ unemployed
    \item $+0.08$ multiple problem debts
    \item $+0.05$ vague about income
\end{itemize}
\noindent
A confidence cap is applied based on description quality (extremely short
text caps at 0.45; rich detail caps at 0.95).

\noindent
\textbf{Role:} first-pass distillation of prose into structured numbers.
Designed for scalability (runs on every loan).

\subsection{Feature Strategist (LLM, White Agent 2)}

\noindent
\textbf{Input:} description text, baseline probability, text risk and
confidence, the gate's per-row weight (or static grade weight),
grade-level training history, and the previous iteration's Advocate
opportunity flags (in the training loop).

\noindent
\textbf{Output:} a \texttt{strategy\_type} (\texttt{keep\_learned},
\texttt{case\_specific\_exception}, \texttt{fuzzy\_only},
\texttt{explanation\_only}, or \texttt{human\_review}), a
\texttt{proposed\_weight}, optional constraints, a \texttt{rationale},
and (in the training notebook) a per-grade signal classification
(\texttt{PERSISTENT\_WIN} / \texttt{PERSISTENT\_LOSS} / \texttt{NOISY} /
\texttt{UNKNOWN}).

\noindent
\textbf{Trigger conditions:}
\begin{itemize}[itemsep=0pt]
    \item In all modes: \texttt{|text\_risk - baseline\_p| > 0.20} (real disagreement)
    \item In static mode only: baseline in fuzzy zone $[0.30, 0.65]$, or extreme text
          with a mismatched grade weight
\end{itemize}
\noindent
In gate mode, the fuzzy-zone trigger is dropped because the gate already
incorporates \texttt{baseline\_p} as a feature.

\noindent
\textbf{Role:} \emph{meta-judgment} about whether the gate's weight needs
case-specific correction based on description-level cues the gate cannot
read (see \S\ref{sec:strategist-vs-analyst} for the detailed analysis).

\subsection{Subgroup Advocate (LLM, White Agent 3)}

\noindent
\textbf{Input:} previous and current fusion result (with per-grade AUC and
AUPRC deltas), Green's recent prediction trajectory, sample sizes per grade.

\noindent
\textbf{Output:} \texttt{veto} (true/false), \texttt{affected\_grades},
\texttt{critique}, \texttt{constraints}, \texttt{opportunity\_flags}, and
\texttt{reasoning}.

\noindent
\textbf{Decision principles (in the prompt):}
\begin{itemize}[itemsep=0pt]
    \item Veto only if harm is \textbf{persistent} (negative delta in $\geq 2$ recent
          iterations) OR a single-step drop clearly exceeds sampling noise.
    \item Account for sample size: small $n$ means high noise, so require a
          larger drop to veto.
    \item Flag \textbf{opportunities} (not vetoes): grades with consistently
          positive deltas worth nudging up.
    \item Be conservative about vetoing; false vetoes block real gains.
\end{itemize}
\noindent
A rule-based fallback fires if the LLM call fails.

\noindent
\textbf{Role:} subgroup safety check. In the demo, only fires when the
Strategist proposes a non-trivial deviation from the baseline weight.

\subsection{Green Agent (hub role)}

Green plays seven distinct functions:

\begin{enumerate}[itemsep=0pt]
    \item \textbf{Executor:} runs the XGBoost baseline, computes fused predictions,
          per-grade AUC \emph{and} AUPRC.
    \item \textbf{Validator:} checks if a proposed strategy passes (mean shift,
          flip-stat precision and recall) with \emph{adaptive thresholds}
          derived from Reflexion experience.
    \item \textbf{Reflexion planner:} before each training iteration, generates
          an \emph{expectation} of $\Delta$AUPRC, mean shift, and per-grade deltas;
          after evaluation, \emph{reflects} on prediction error and accumulates
          \texttt{self.experience}.
    \item \textbf{Arbitrator:} on Advocate veto, synthesizes the Strategist's
          proposal, the Advocate's constraints, and Green's own experience
          into a compromise weight. (See \S\ref{sec:design-arbitrator}.)
    \item \textbf{Terminal decider:} \texttt{green\_llm\_review} when all
          candidates fail validation.
    \item \textbf{Gate trainer:} learns the Tier-A per-row text-weight function
          via scipy L-BFGS-B (see \S\ref{sec:gate}).
    \item \textbf{Final authority:} flip-gate hard rule before the verdict is
          issued (verdict reversals require high confidence + extreme text signal).
\end{enumerate}
\noindent
This concentration of roles is intentional: many of them share the same
\texttt{self.experience} memory and \texttt{evaluate\_fusion} computation, so
splitting them would create awkward state-sharing.

\subsection{Reporter (LLM, White Agent 4)}

\noindent
\textbf{Input:} the entire context (baseline, fused, weights, grade, text
risk, confidence, the borrower's key numeric features, the decision threshold).

\noindent
\textbf{Output:} two short paragraphs in natural language describing
(1) the numeric-risk picture and (2) the text's role in shifting the prediction.

\noindent
\textbf{Role:} human-readable summary. Reads only after fusion completes;
never affects the prediction.

%==============================================================
\section{The Tier-A learned gating function}
\label{sec:gate}

\subsection{Motivation}

The original system used five hand-tuned scalars
\texttt{grade\_weights = \{C: 0.03, D: 0.13, E: 0.08, F: 0.06, G: 0.03\}}
to weight text influence per grade. This is dumb in two ways:
(1) it does not vary within a grade, even though some cases within Grade C
might benefit greatly from text while others should not; (2) it cannot
respond to text content (a protective borrower's note gets the same weight
as an alarming one if they share a grade).

The Tier-A gate replaces these five scalars with a learned per-row function:
\[
w_{\mathrm{text}}(\mathrm{row}) = w_{\max} \cdot \sigma\bigl(\mathbf{x} \cdot \mathbf{w} + b\bigr)
\]
where $\mathbf{x}$ is an 11-dimensional feature vector and the parameters
$(\mathbf{w}, b)$ are optimized to minimize the log-loss of the fused
prediction on the training set.

\subsection{Input features}

The gate sees:
\begin{itemize}[itemsep=0pt]
    \item Five grade one-hot indicators (C, D, E, F, G)
    \item \texttt{baseline\_p} (XGBoost output)
    \item \texttt{text\_conf} (Text Analyst's confidence)
    \item \texttt{text\_risk} (Text Analyst's score)
    \item Three derived features: $|\texttt{text\_risk} - 0.5|$ (text extremity),
          $|\texttt{baseline\_p} - 0.5|$ (numeric uncertainty),
          $|\texttt{text\_risk} - \texttt{baseline\_p}|$ (model disagreement)
\end{itemize}
\noindent
Crucially, \textbf{the gate does not see the description prose}. It sees only
the numeric distillations produced by Text Analyst.

\subsection{Training objective}

Let $y_i$ be the binary default label, $b_i$ the baseline probability, $t_i$
the text risk score, $c_i$ the text confidence, and $\mathbf{x}_i$ the
standardized feature vector. Then for parameters $(\mathbf{w}, b)$:
\[
w_i = w_{\max} \cdot \sigma(\mathbf{x}_i \cdot \mathbf{w} + b) \quad\quad
\hat{p}_i = \sigma\bigl(\mathrm{logit}(b_i) + w_i \, c_i \, (\mathrm{logit}(t_i) - \mathrm{logit}(0.5))\bigr)
\]
\[
\mathcal{L}(\mathbf{w}, b) = -\frac{1}{N}\sum_{i=1}^{N} \bigl[y_i \log \hat{p}_i + (1-y_i)\log(1-\hat{p}_i)\bigr] + \lambda \|\mathbf{w}\|^2
\]
Optimized using L-BFGS-B from scipy with $\lambda = 10^{-3}$. The fusion
formula is differentiable end-to-end through $w_i$.

\subsection{What the gate learned}

After training (on cached \texttt{subset\_train} with $n = 880$), the dominant
coefficients are:

\begin{center}
\begin{tabular}{lr}
\toprule
\textbf{Feature (post-standardization)} & \textbf{Coefficient} \\
\midrule
\texttt{text\_risk} & $-2.15$ \\
\texttt{|text\_risk - 0.5|} & $-0.40$ \\
\texttt{|text\_risk - baseline\_p|} & $+0.35$ \\
\texttt{text\_conf} & $+0.34$ \\
\texttt{baseline\_p} & $+0.13$ \\
bias & $+0.36$ \\
\bottomrule
\end{tabular}
\end{center}
\noindent
The strongly negative \texttt{text\_risk} coefficient means: the lower the
text risk score (i.e., the more protective the text), the higher the
gate-output weight. The gate has learned the asymmetry that
\emph{protective text deserves to be trusted more strongly than risky text}
(the latter is likely already reflected in the numeric baseline).

\subsection{Persistence}

After training, the gate parameters $(\mathbf{w}, b, \boldsymbol{\mu},
\boldsymbol{\sigma}, \texttt{grades}, w_{\max})$ are serialized as plain
Python lists/floats into the model pkl under key
\texttt{gating\_params}. The Streamlit demo loads this and reconstructs
the per-row inference at run time. If the key is absent the demo falls
back to the static 5-scalar policy.

%==============================================================
\section{Worked example}
\label{sec:example}

\textbf{Input:} A Grade C borrower with the following numeric profile and
description:

\begin{tcolorbox}[colback=codebg,colframe=gray!50,boxrule=0.5pt]
\textit{Loan}: \$14k at 17.5\%, FICO 705--709, DTI 14\%, revol\_util 38\%, \ldots

\textit{Description}: ``I am a registered nurse with 9 years of continuous
employment at the same hospital. My annual salary is \$68,000. I have never
missed a single payment on any account. This loan consolidates two low-balance
credit cards.''
\end{tcolorbox}

\subsection{Stage trace}

\begin{tabular}{@{}p{0.05\linewidth}p{0.20\linewidth}p{0.40\linewidth}p{0.25\linewidth}@{}}
\toprule
\textbf{\#} & \textbf{Stage} & \textbf{Computation} & \textbf{Output} \\
\midrule
1 & Green Baseline & XGBoost on 16 numeric features & $p_{\mathrm{base}} = 0.624$ (DEFAULT zone) \\
\addlinespace
2 & Text Analyst & LLM scores description; $0.45 - 0.15 - 0.10 - 0.05 \to 0.15$, clamp $\to 0.05$ & $\texttt{text\_risk}=0.05, \texttt{conf}=0.90$ \\
\addlinespace
3 & Tier-A Gate & Apply learned function: dominant $\texttt{text\_risk}=-2.15$ coefficient $\times$ low risk $\to$ high weight & $w_{\mathrm{gate}} = 0.598$ \\
\addlinespace
4 & Strategist & Triggered (score gap $0.574 > 0.20$). LLM reads prose, sees gate already captured everything. Decision framework: ``keep\_learned is the right default.'' & \texttt{keep\_learned}, $w=0.598$ \\
\addlinespace
5 & Advocate & \textit{Skipped} --- Strategist proposed no change & --- \\
\addlinespace
6 & Green Arbitrate & \textit{Skipped} --- Advocate did not veto & --- \\
\addlinespace
7 & Fusion & $\mathrm{logit}(0.624) + 0.598 \cdot 0.90 \cdot (\mathrm{logit}(0.05) - 0) = -1.082$ & \texttt{fused\_pred} $= 0.253$ \\
\addlinespace
8 & Green Final Authority & Shift cap \textbf{disabled} (gate mode). Flip gate detects risky$\to$safe; requires $\texttt{conf} \geq 0.87$ (got 0.90 \checkmark) and $|\texttt{text}-0.5| \geq 0.35$ (got 0.45 \checkmark). Flip approved. & $\texttt{green\_action}=\texttt{accepted}$ \\
\addlinespace
9 & Green LLM Review & \textit{Skipped} --- no hard rule fired & --- \\
\addlinespace
10 & Reporter & LLM generates explanation: text shifted probability from 0.624 to 0.253, changing verdict to NON-DEFAULT & 2 paragraphs \\
\bottomrule
\end{tabular}

\subsection{Final verdict: NON-DEFAULT (0.253)}

This example illustrates the system's correct behavior: protective text
with high confidence \emph{does} get to flip a numeric verdict, but only
when both the gate's structural signal and Green's outcome-level safeguards
agree.

\subsection{Which agents actually fired?}

\begin{center}
\begin{tabular}{@{}lcl@{}}
\toprule
\textbf{Agent} & \textbf{Fired?} & \textbf{Role on this case} \\
\midrule
Green Baseline & \checkmark always & gave 0.624 \\
Text Analyst & \checkmark always (text present) & prose $\to$ (0.05, 0.90) \\
\textbf{Tier-A Gate} & \checkmark always & \textbf{the actual decision-maker:} 0.598 \\
Strategist & \checkmark (score gap) & \textbf{chose not to intervene} --- judgment \\
Advocate & --- skipped & (no change to veto) \\
Green Arbitrate & --- skipped & (no veto) \\
Fusion & \checkmark always & 0.253 \\
Green Final & \checkmark always & flip gate approved \\
Green LLM Review & --- skipped & (no rule triggered) \\
Reporter & \checkmark always & explanation \\
\bottomrule
\end{tabular}
\end{center}

\noindent
\textbf{The key observation:} most cases do \emph{not} require Strategist,
Advocate, Arbitrate, or LLM Review. The Tier-A gate already produces
per-case weights that work for the typical loan. The LLM-based
conditional agents fire only when there is genuine signal disagreement
or rule-level contention.

%==============================================================
\section{Design discussion: Strategist vs.\ Text Analyst overlap}
\label{sec:strategist-vs-analyst}

\subsection{The concern}

A natural question: \emph{both the Text Analyst and the Strategist read the
description. Doesn't the Strategist duplicate the Analyst's work?}

The gate sees only the Analyst's compressed output
\texttt{(text\_risk, text\_conf)} and never the original prose. The
Strategist sees both --- so if it is just re-scoring the same risk
signals, it is in fact redundant.

\subsection{What the Text Analyst loses by compressing}

The Analyst follows a rule-based rubric (\S\ref{sec:agents}.1). This
compression necessarily loses five categories of signal:

\begin{center}
\begin{tabular}{@{}p{0.22\linewidth}p{0.36\linewidth}p{0.36\linewidth}@{}}
\toprule
\textbf{Signal type} & \textbf{Example} & \textbf{Analyst's handling} \\
\midrule
Verifiability & ``Microsoft engineer 8 yrs'' vs.\ ``I have a good job'' & Both trigger ``$-0.15$ stable employment'' \\
\addlinespace
Internal contradiction & ``Stable job 9 yrs, but currently behind on rent'' & Mixed signals $\to$ middle score with low confidence \\
\addlinespace
Detail specificity & ``\$68K salary'' vs.\ ``decent income'' & First triggers $-0.08$, second does not, but confidence may not reflect the gap \\
\addlinespace
Life-event context & ``divorce + stable income'' vs.\ ``new job + 30\% raise'' & Both might produce similar middle scores \\
\addlinespace
``Too good to be true'' & baseline$=0.7$ (high risk) plus a perfect-record text & Analyst never sees the baseline, so cannot question consistency \\
\bottomrule
\end{tabular}
\end{center}

\noindent
The last row is the most architecturally significant: the Text Analyst
operates in isolation from the numeric model. By construction, it cannot
ask ``does this text contradict what the numbers say?''

\subsection{What the Strategist sees that nobody else does}

\begin{center}
\begin{tabular}{@{}lcc@{}}
\toprule
& \textbf{Text Analyst} & \textbf{Strategist} \\
\midrule
description prose & \checkmark & \checkmark \\
\texttt{baseline\_p} (XGBoost output) & --- & \checkmark \\
\texttt{text\_risk}, \texttt{text\_conf} & (own output) & \checkmark \\
gate's per-row weight & --- & \checkmark \\
grade-level training history & --- & \checkmark \\
\bottomrule
\end{tabular}
\end{center}

\noindent
The Strategist's task type is also different: \emph{first-order} judgment
(``what does the text say?'') vs.\ \emph{second-order} meta-judgment
(``how much should we trust it given the case context?'').

\subsection{Acknowledged risk of duplication}

If the Strategist's prompt is poorly designed, it can collapse into a
re-do of the Text Analyst. The current Strategist prompt counters this with
an explicit decision framework:
\begin{quote}\itshape
``Read the borrower's description prose. Does it contain CASE-SPECIFIC
evidence --- employer reputation, repayment-plan specifics, life-event
context, urgency cues, contradictions with the numeric profile --- that
the gate's structured features cannot see?''
\end{quote}
\noindent
The framework explicitly anchors the LLM's role to ``what the gate cannot
see.'' In gate mode, the default outcome is \texttt{keep\_learned}; the
Strategist must \emph{actively justify} any deviation by citing a specific
phrase from the description.

\subsection{Validation test}

A quick empirical check for duplication:
\begin{enumerate}[itemsep=0pt]
    \item \textbf{Generic-signal case:} a description containing only signals the
          Analyst already covers. The Strategist should choose
          \texttt{keep\_learned}; if it proposes a deviation, it is being
          redundant.
    \item \textbf{Contradiction case:} ``stable job 9 yrs'' \emph{and} ``behind on
          rent.'' The Strategist should propose \texttt{case\_specific\_exception}
          with the contradiction cited explicitly. If it just repeats the
          Analyst's risk/protective list, it is duplicating.
\end{enumerate}

\subsection{Possible refinements (future work)}

\begin{enumerate}[itemsep=0pt]
    \item Hide \texttt{(text\_risk, text\_conf)} from the Strategist --- force
          it to read prose only. \emph{Risk}: now both agents do raw-text
          scoring, the real duplication.
    \item Hide prose from the Strategist --- give it only the structured numbers.
          \emph{Risk}: it loses the case-specific detail that justifies its existence.
    \item Restrict Strategist to a ``contradiction auditor'' role --- only fires
          on internal inconsistencies the Analyst could not flag. \emph{Trade-off}:
          cleanest division of labor, but limits the Strategist's scope.
\end{enumerate}

%==============================================================
\section{Evaluation methodology}
\label{sec:eval}

The evaluation pipeline answers six concrete questions, each backed by a
metric chosen to match the question, not by convention.

\subsection{AUPRC as primary discrimination metric}

AUC is the conventional choice but is misleading at imbalanced base
rates (loan default is $\sim 15$--$25\%$ in C--G grades). AUPRC weights
the positive class more honestly. All loop prints, decision logs, KPI
milestones, Tier-A comparisons, and final test evaluations now lead with
AUPRC; AUC is shown in parentheses as a secondary diagnostic.

\subsection{The six-question toolkit}

\begin{center}
\begin{tabular}{@{}p{0.45\linewidth}p{0.50\linewidth}@{}}
\toprule
\textbf{Question} & \textbf{Metric} \\
\midrule
Does fusion improve discrimination at this base rate? & AUPRC with bootstrap CI \\
Does fusion break calibration? & Brier score + ECE + per-bin reliability \\
Across plausible decision thresholds, is fusion net-beneficial? & Decision-curve net benefit \\
When fusion changes the verdict, does it reclassify correctly? & Categorical and continuous NRI \\
Is the LLM's self-reported confidence informative? & LLM confidence reliability plot \\
\textbf{When should we trust the text?} & Selective fusion curve \\
\bottomrule
\end{tabular}
\end{center}

\noindent
The \textbf{selective fusion curve} is the headline measurement: for each
candidate ``trust score'' (text confidence, model disagreement, or the
Tier-A gate's own output), we sort cases by that score and plot the
fused $\Delta$AUPRC as we sweep coverage from top-5\% down to 100\%.
The curve's shape \emph{is} the answer to ``when to trust the text'':
monotonically decreasing means trust only at the top, inverted-U means
there is an optimal coverage threshold, and a flat curve means the
trust signal carries no useful information.

\subsection{Per-Agent KPI (MultiAgentBench \S3.3)}

For per-agent attribution we adopt the milestone-based KPI from
MultiAgentBench (Zhu et al., 2025):
\[
\mathrm{KPI}_j = \frac{n_j}{M}, \quad \mathrm{KPI}_{\mathrm{overall}} = \frac{1}{N}\sum_{j=1}^{N}\mathrm{KPI}_j
\]
where $n_j$ is the number of milestones agent $j$ contributed to and $M$
is the total milestone count. We define 11 milestones --- 7 per-iteration
and 4 once-per-run --- with explicit (deterministic, non-LLM) attribution
rules. Each milestone names the contributing agents.

\noindent
This surfaces the system's \emph{hub} structure: Green typically scores
the highest KPI because it covers 7 distinct roles. Strategist and
Advocate contribute to fewer milestones but in load-bearing places
(the strategy-proposed and no-veto milestones, respectively).

%==============================================================
\section{Architectural decisions}
\label{sec:decisions}

\subsection{Arbitrator merged into Green}
\label{sec:design-arbitrator}

The original architecture had Arbitrator as a separate White Agent. We
merged it into Green for three reasons:

\begin{enumerate}[itemsep=0pt]
    \item Arbitrator's job (mediate Strategist vs.\ Advocate conflict) requires
          access to the same Reflexion experience and \texttt{evaluate\_fusion}
          machinery that Green already owns.
    \item Splitting it into a separate agent forced redundant state sharing
          and made the trajectory of past arbitrations awkward to retrieve.
    \item Conceptually, arbitration is an authority decision, and Green is
          the system's authority agent (final flip gate, terminal LLM review,
          baseline owner). All three roles are now consistent.
\end{enumerate}
\noindent
The LLM call and the prompt are essentially unchanged; only the
\emph{identity} attribution moved.

\subsection{Decision threshold unified to 0.5}

A prior version of the system had a calibrated threshold ($0.62$ in the
training run) computed by \texttt{find\_threshold} to enforce a
precision-$\geq 0.40$ floor. This was inconsistent with every other part
of the system, all of which used $0.5$ (per-grade precision/recall, flip
stats, evaluation, etc.).

The calibrated threshold is now \emph{informational only}: we still
report ``if you want precision $\geq 0.40$, threshold would need to be
$0.62$,'' but the actual decision threshold is fixed at $0.5$. This
removes a confusing inconsistency without losing the diagnostic.

\subsection{MAX\_SHIFT cap disabled in gate mode}

In static mode (5-scalar weights), Green enforces a maximum probability
shift of $0.12$ from baseline: a safety rail against text dominating
the numeric model. This made sense when weights were dumb constants.

In gate mode, the gate has explicitly learned \emph{when a large shift
is appropriate}. Capping the shift at 0.12 directly contradicts what
the gate learned. The shift cap is now disabled when
\texttt{gating\_params} is present in the pkl. The asymmetric
flip-gate hard rule (verdict reversals require both high confidence
and extreme text) remains active --- but operates on outcomes rather
than on weights, so it does not fight the gate.

\subsection{Reflexion in Green}

Green maintains \texttt{self.experience} across loop iterations. Before
each evaluation, \texttt{generate\_expectation} predicts the
$\Delta$AUPRC, mean shift, and per-grade deltas. After evaluation,
\texttt{reflect} records the actual values and computes the prediction
error. The accumulated history drives:

\begin{itemize}[itemsep=0pt]
    \item \textbf{Adaptive validation thresholds:} max-mean-shift is derived
          from observed shift std rather than a fixed constant.
    \item \textbf{Memory-aware fallback:} \texttt{propose\_counter\_strategy}
          uses recent per-grade trends to adjust weights when LLM arbitration
          fails.
    \item \textbf{Trajectory injection into prompts:} both
          \texttt{green\_llm\_review} and \texttt{arbitrate} include the
          recent prediction trajectory as context, so the LLM can reason
          about systematic over/under-estimation.
\end{itemize}

%==============================================================
\section{Open questions and future work}
\label{sec:future}

\subsection{Strategist refinement}

The current Strategist is gate-aware and chooses \texttt{keep\_learned} as
the default, but its prompt does not formally enforce ``find a contradiction
or specific phrase, otherwise keep.'' The \texttt{contradiction-auditor}
refinement (\S\ref{sec:strategist-vs-analyst}.6) would make this clearer.

\subsection{Text Analyst depth}

The current Text Analyst uses a rule-based scoring rubric. A richer analyst
that extracts structured fields (employer name, repayment plan specificity,
internal contradictions) could reduce the Strategist's overlap concern by
shifting more of the prose interpretation upstream.

\subsection{Tier-A overfitting}

In the executed run, the Tier-A gate's training log-loss improved by
$+0.030$, but on the small test set ($n \approx 80$ cached observations)
the gate's test AUPRC ($0.39$) underperformed the static Strategist
weights ($0.43$). The gate has 11 parameters and was trained on
$\sim 880$ rows --- borderline but not pathological. A larger Text
Analyst cache would let us re-evaluate whether the gap is overfitting
or test-set noise.

\subsection{Cost control}

The current configuration makes 5--6 LLM calls per training iteration
(Strategist, Advocate, Green expectation, possibly Green arbitrate,
possibly Green LLM review). With \texttt{max\_iter}$=5$ this is
$25$--$30$ calls per training run plus walkthroughs. An optional
\texttt{llm\_intensity} flag (\texttt{low} / \texttt{medium} / \texttt{high})
would let users trade off cost and decision quality.

\subsection{Topology experiments}

MultiAgentBench reports that graph-mesh topologies outperform chain and
tree in research-style tasks. The current pipeline is essentially a chain
(Strategist $\to$ Advocate $\to$ Green). A graph-mesh variant --- where
each White Agent independently sees the others' outputs in parallel ---
would be an informative ablation, especially given Green's existing hub
role.

%==============================================================
\section{Conclusion}

The system is a five-agent collaboration on loan-default prediction with
a learned per-row gating function (Tier-A) at its core. The Green Agent
owns the multiple authority roles --- baseline, gate, arbitration, flip
gate, terminal review --- and uses Reflexion memory to adapt its
validation thresholds and arbitration prompts across iterations.

The two architectural insights this document tries to convey:

\begin{enumerate}[itemsep=0pt]
    \item \textbf{Most cases do not require LLM intervention beyond the
    Text Analyst and the Tier-A gate.} The Strategist, Advocate,
    Arbitrate, and LLM Review agents fire only when the gate's
    structural signal might be inadequate. The system's frugality with
    LLM calls is a feature, not a sign of underuse.

    \item \textbf{The Strategist's value is meta-judgment, not re-scoring.}
    The Text Analyst tells us what the text says; the Strategist tells us
    whether to trust that signal given the case's full context. Whether
    this distinction holds up in practice is the subject of ongoing
    refinement (\S\ref{sec:strategist-vs-analyst}).
\end{enumerate}

\end{document}
